\section{周王室之二：成、康、穆之间的“安定”从哪里来}

后世常把成王、康王时期概括为“成康之治”。这个词有它的价值：它提醒我们，周初在东征、营成周和封国布局之后，确实有一段王室权威较能运转的时期。但它也容易把几十年的政治、军事、赏赐和地方经营压缩成一个没有摩擦的黄金时代。

\begin{event}{成王、康王时期}{东征后的秩序巩固}{相对次序较可靠}{宗周、成周与东方封国}\eventanchor{WZ-CHENGKANG-CONSOLIDATION}{西周·成康巩固}
\term{这一阶段发生了什么}：周初危机平定后，王室继续通过册命、赏赐、征伐和祭祀维持关系网络。何尊、大盂鼎等铭文所保留的，不是平静无事的“盛世画面”，而是一套不断确认功劳、授予权利、要求服从的政治日常。

\term{事情为什么重要}：新制度能活下来，不靠一次立法，而靠一代人反复把它做成习惯。宗周和成周的双重布局、东方封国的经营、殷遗民的安置，都需要在多年里不断投入资源。

\term{后来改变了什么}：王室的安定为后续向南、向北的行动提供了余地；也让更大规模的封国网络逐渐成为既成事实。网络越成熟，中心越需要持续支付维系它的成本。
\end{event}

\begin{event}{穆王时期}{王室远行与后出的传奇}{材料层次不一}{西北与诸侯地区}\eventanchor{WZ-MU-TRAVELS}{西周·穆王远行}
\term{这一阶段发生了什么}：传世材料把周穆王描绘成频繁巡行、征伐和西行的君主。部分西周中期金文可见王室军事与册命活动；但《穆天子传》式的远行故事包含浓厚的后世文学和想象成分。

\term{该怎样读}：穆王故事有助于理解后人如何想象一个强盛王权的活动范围，却不能被直接当成地理考察报告。可靠的历史问题是：王室为何需要不断巡行、征伐和赏赐？答案仍在于它必须让远方政治体不断感到中心的存在。
\end{event}

\begin{textwindow}[“穆天子”不是旅行攻略]
《穆天子传》中的西行叙事极富魅力，但它的成书、文本层次和历史真实性都需要谨慎判断。把它放在这里，是为了让读者看见西周王权在后世记忆中如何被想象为可以走到遥远边疆的力量，而不是把传说中的路线直接画进历史地图。
\end{textwindow}
